\begin{center}
\large{\MakeUppercase{\textbf{Анотація}}}
\end{center}

\newcommand{\keyword}{Ключові слова:}

Пояснювальна записка містить 69 сторінок, 10 рисунків, 2 таблиці, 38 джерел.

У даній дипломній роботі проведено огляд та аналіз існуючих мобільних застосунків для відстеження розвитку немовляти, досліджено їх переваги та недоліки. На основі аналізу предметної області та потреб користувачів сформульовано вимоги до системи, обрано технологічний стек і спроектовано архітектуру кросплатформного мобільного застосунку. Розроблено реляційну схему бази даних із захистом на рівні рядків (Row-Level Security), що забезпечує ізоляцію даних між користувачами, та визначено ключові модулі системи.

В результаті розроблено кросплатформний мобільний застосунок для платформ iOS та Android, що поєднує щоденник активностей немовляти, моніторинг антропометричних показників відповідно до стандартів ВООЗ, управління вакцинаціями та досягненнями, а також ШІ-асистента на основі технології Retrieval-Augmented Generation. Застосунок реалізовано з використанням React Native та Expo; серверну частину побудовано на Supabase із PostgreSQL, pgvector та Row-Level Security. ШІ-асистента інтегровано через Groq API з векторним пошуком у базі педіатричних знань. Розгортання здійснено на хмарній інфраструктурі з підтримкою синхронізації між пристроями в реальному часі. Проведено тестування системи на рівнях юніт-тестів, тестів бази даних і наскрізних тестів.

\keyword\ мобільний застосунок, React Native, Expo, Supabase, pgvector, RAG, ШІ-асистент, відстеження розвитку немовляти, ВООЗ.


\newpage

\begin{center}
\large{\MakeUppercase{\textbf{Annotation}}}
\end{center}

The explanatory note contains 69 pages, 10 figures, 2 tables, 38 references.

In this bachelor's thesis, a review and analysis of existing mobile applications for infant development tracking was conducted, and their advantages and shortcomings were examined. Based on the subject domain analysis and user needs, system requirements were formulated, a technology stack was selected, and the architecture of a cross-platform mobile application was designed. A relational database schema with Row-Level Security was developed to ensure data isolation between users, and the key system modules were defined.

As a result, a cross-platform mobile application for iOS and Android platforms was developed, combining an infant activity journal, monitoring of anthropometric indicators in accordance with WHO standards, management of vaccinations and developmental milestones, and an AI assistant based on Retrieval-Augmented Generation technology. The application was implemented using React Native and Expo; the backend was built on Supabase with PostgreSQL, pgvector and Row-Level Security. The AI assistant was integrated via the Groq API with vector search in a paediatric knowledge base. The system was deployed on cloud infrastructure with support for real-time synchronisation between devices. System testing was conducted at unit, database and end-to-end levels.

Keywords: mobile application, React Native, Expo, Supabase, pgvector, RAG, AI assistant, infant development tracking, WHO.

\newpage
